\chapter{Bradycardia}

\section*{Definition}
Bradycardia is defined as a resting heart rate less than 60 beats per minute in adults. It is considered pathologic when accompanied by symptoms (syncope, presyncope, dyspnea, fatigue) or when the ventricular rate is inadequate to maintain cardiac output.

\section*{Pathophysiology}
Bradycardia results from decreased automaticity of the sinoatrial node, conduction block through the atrioventricular node, or activation of vagal tone. The reduced heart rate decreases cardiac output (heart rate $\times$ stroke volume), which may compromise cerebral and systemic perfusion.

\section*{Etiology}
\begin{itemize}
  \item \textbf{Physiologic:} Endurance athletes (increased vagal tone, resting HR 30--50), sleep, trained cardiovascular conditioning
  \item \textbf{Cardiac:} Sick sinus syndrome, AV block (first, second, third degree), ischemic heart disease, myocarditis, infiltrative disease (sarcoidosis, amyloidosis)
  \item \textbf{Pharmacologic:} Beta-blockers, calcium channel blockers (verapamil, diltiazem), digoxin, amiodarone, clonidine, ivabradine
  \item \textbf{Metabolic:} Hypothyroidism, hypothermia, hyperkalemia, hypoglycemia, hypocalcemia, severe hypoxemia
  \item \textbf{Neurologic:} Increased intracranial pressure (Cushing reflex), vasovagal reflex, neurocardiogenic syncope
  \item \textbf{Other:} Obstructive sleep apnea, anorexia nervosa, Lyme carditis, COVID-19, post-cardiac surgery
\end{itemize}

\section*{Indications for Evaluation}
Seek care if:
\begin{itemize}
  \item Syncope, presyncope, or near-fainting episodes
  \item Chest pain, dyspnea, or altered mental status with bradycardia
  \item Known heart disease, pacemaker, or recent cardiac surgery
  \item Heart rate less than 40 beats/min with symptoms
  \item Suspected drug toxicity (digoxin, beta-blocker overdose)
\end{itemize}

\section*{Related Symptoms}
\begin{itemize}
  \item Fatigue, dizziness, lightheadedness, syncope, dyspnea, exercise intolerance, chest pain, confusion
\end{itemize}
\textbf{Note:} Asymptomatic bradycardia (especially in athletes or during sleep) generally requires no intervention; the decision to treat is based on symptoms, not the absolute heart rate number alone.

\section*{Self-Care / Home Management}
\begin{itemize}
  \item Identify and discontinue offending medications or substances under medical supervision
  \item Avoid triggers: Valsalva, straining during bowel movements, rapid postural changes
  \item Increase oral hydration if volume depletion contributes to vagal tone
  \item Monitor heart rate and blood pressure at home; log symptomatic episodes
  \item Report any syncopal episodes or significant changes to the treating clinician
\end{itemize}

\section*{Diagnostic Approach}
\begin{itemize}
  \item 12-lead ECG to identify rhythm: sinus bradycardia, AV block (PR interval, Mobitz I/II, complete heart block), junctional or idioventricular escape rhythm
  \item Holter monitor (24--48 hours) or extended cardiac event monitor for intermittent symptoms
  \item Exercise stress testing to assess chronotropic competence
  \item Echocardiogram to evaluate structural heart disease and left ventricular function
  \item Laboratory: thyroid function tests, electrolytes (especially potassium, calcium), cardiac troponin, digoxin level if applicable
\end{itemize}

\section*{Treatment / Management Overview}
\begin{itemize}
  \item Asymptomatic physiologic bradycardia: no treatment required
  \item Withhold or adjust dose of causative medications
  \item Acute symptomatic bradycardia: IV atropine (0.5 mg every 3--5 min, max 3 mg) as first-line; transcutaneous pacing for refractory cases
  \item Epinephrine or dopamine infusion for hemodynamically unstable bradycardia
  \item Permanent pacemaker insertion for symptomatic sinus node dysfunction or high-grade AV block (Mobitz II, complete heart block)
\end{itemize}

\clearpage
